\chapter{A complete search by hand}
\label{ch:byhand}

\begin{objectives}
\begin{itemize}
\item Execute eight full iterations of Leela's search with pencil and paper --- selection,
  expansion, evaluation, backpropagation --- and match the engine's own numbers to five
  decimal places.
\item Watch the search deepen along its main line without being told to.
\item Watch a refutation happen: a move examined once, scored $0.000$, and abandoned ---
  with independent confirmation that abandoning it was correct.
\item Explain which number finally chooses the move played, and why it is not $Q$.
\item Discover, by experiment, that evaluating a position from a bare FEN is not the same as
  evaluating it in a game.
\end{itemize}
\end{objectives}

\section{The position, and the truth about it}

Everything in this chapter is real. The priors, the values, the selections and the resulting
tree are what \texttt{engine/lc0.exe} with \texttt{791556.pb.gz} actually did on this
machine; the arithmetic below is reproduced by \texttt{tools/simulate\_search.py}, which
checks itself against the engine and refuses to disagree with it.

\begin{center}
\chessboard[setfen=4k3/8/4K3/4P3/8/8/8/8 w - - 0 1, boardfontsize=19pt]
\end{center}

\noindent White: Ke6, pawn e5. Black: Ke8. \textbf{White to move.} Four legal moves: Kd6,
Kf6, Kd5, Kf5 (the pawn cannot advance --- its own king is in front of it, and d7, e7, f7 are
all controlled by the black king).

We chose this position for one reason: with four legal moves the whole tree fits on a page.
But it is not a toy --- it contains a real trap, and here is the independent ground truth,
from Stockfish 16.1 at depth 30 (\emph{not} from Leela, and not from the author):

\begin{center}
\begin{tabular}{@{}llp{7.3cm}@{}}
\toprule
Move & Stockfish depth 30 & \\
\midrule
\textbf{Kd6} & \textbf{mate in 12} & wins \\
\textbf{Kf6} & \textbf{mate in 18} & also wins, less efficiently \\
Kd5 & $0.00$ & \textbf{throws the win away} \\
Kf5 & $0.00$ & \textbf{throws the win away} \\
\bottomrule
\end{tabular}
\end{center}

So two of the four moves win and two draw. Keep that in your pocket; the search does not know
it yet.

\section{The bookkeeping}

Every iteration we recompute, for each of the four moves,
\[
  S(a) \;=\; \underbrace{Q(a)}_{\text{measured, or FPU if unmeasured}}
  \;+\;\underbrace{c_{\text{puct}}\,P(a)\frac{\sqrt{N}}{1+n_a}}_{U(a)}
\]
and select the largest. The constants, from Chapter~\ref{ch:puct}:
\[
  c_{\text{puct}}(N) = 1.745 + 3.894\ln\!\frac{N + 38740}{38739},\qquad
  Q_{\text{FPU}} = Q(\text{root}) - 0.33\sqrt{\textstyle\sum_{\text{visited}} P},
\]
and $N$ means the total visits already spent on the root's children (using $\sqrt{\max(N,1)}$
so that the first selection is decided by the priors alone).

The network's opening statement about this position --- one forward pass, no search:

\begin{realdata}[title={Real engine output: the root evaluation and priors}]
\[
  V(\text{root}) = +0.97602, \qquad d = 0.024
\]
\begin{center}
\begin{tabular}{@{}lrl@{}}
\toprule
Move & $P(a\given s)$ & \\
\midrule
Kd6 & 45.13\% & \\
Kf6 & 44.23\% & \\
Kf5 &  5.38\% & one of the two moves that draws \\
Kd5 &  5.26\% & the other \\
\bottomrule
\end{tabular}
\end{center}
Note what the policy has and has not done. It has correctly identified that the two
king-forward moves are the serious ones --- $89\%$ of its attention goes there. But it still
assigns $10.6\%$ to two moves that throw away a won game. \textbf{The policy is a good guess,
not an oracle}; the search's job in this position is to spend a little of its budget
converting that $10.6\%$ of doubt into knowledge.
\end{realdata}

\section{The eight iterations}

Do these with a calculator. Each table is the state \emph{before} the selection.

\subsection*{Iteration 1 --- nothing measured}

$N = 0$ (so $\sqrt{\max(N,1)} = 1$), no move has been visited, so all four get
$Q_{\text{FPU}} = Q(\text{root}) = 0.97602$ (nothing visited, so the penalty
$0.33\sqrt{0} = 0$).

\[
  U(\text{Kd6}) = 1.7451 \times 0.4513 \times \frac{1}{1+0} = 0.78756 .
\]

\begin{center}
\begin{tabular}{@{}lrrrr@{}}
\toprule
Move & $n_a$ & $Q$ & $U$ & $S$ \\
\midrule
\textbf{Kd6} & 0 & $0.97602$ & $0.78756$ & $\mathbf{1.76358}$ \\
Kf6 & 0 & $0.97602$ & $0.77186$ & $1.74788$ \\
Kf5 & 0 & $0.97602$ & $0.09389$ & $1.06991$ \\
Kd5 & 0 & $0.97602$ & $0.09179$ & $1.06781$ \\
\bottomrule
\end{tabular}
\end{center}

With every $Q$ identical, the ranking is purely the priors --- exactly as
Chapter~\ref{ch:puct} predicted. \textbf{Select Kd6.} Expand it, ask the network:
the child (Black to move) evaluates to $-0.96766$ from Black's point of view, so the root
edge Kd6 receives $+0.96766$.

\subsection*{Iteration 2 --- the first measurement, and the FPU penalty appears}

Now one move is measured and three are not. The root's own $Q$ has been updated to
$(0.97602 + 0.96766)/2 = 0.97184$, and the visited prior mass is $0.4513$, so
\[
Q_{\text{FPU}} = 0.97184 - 0.33\sqrt{0.4513} = 0.97184 - 0.22169 = 0.75015 .
\]

\begin{center}
\begin{tabular}{@{}lrrrrl@{}}
\toprule
Move & $n_a$ & $Q$ & $U$ & $S$ & \\
\midrule
\textbf{Kf6} & 0 & $0.75015$ & $0.77190$ & $\mathbf{1.52205}$ & FPU \\
Kd6 & 1 & $0.96766$ & $0.39380$ & $1.36146$ & measured \\
Kf5 & 0 & $0.75015$ & $0.09389$ & $0.84404$ & FPU \\
Kd5 & 0 & $0.75015$ & $0.09180$ & $0.84195$ & FPU \\
\bottomrule
\end{tabular}
\end{center}

\begin{keyidea}
Kd6 has the \emph{best} measured value ($0.96766$ against a mere FPU guess of $0.75015$) and
is \emph{not} selected. Its $U$ halved the moment it was visited --- the denominator went
from $1+0$ to $1+1$. This is the anti-greedy mechanism of Chapter~\ref{ch:manymachines}
doing its work: measuring something makes it less urgent to measure again.
\end{keyidea}

\textbf{Select Kf6.} Expand: value $+0.98598$ to the root edge.

\subsection*{Iteration 3 --- the first genuine contest}

Two measured moves, both excellent. $Q(\text{root}) = 0.97655$; visited prior mass
$0.4513+0.4423 = 0.8936$, so $Q_{\text{FPU}} = 0.97655 - 0.33(0.94530) = 0.66460$.

\begin{center}
\begin{tabular}{@{}lrrrr@{}}
\toprule
Move & $n_a$ & $Q$ & $U$ & $S$ \\
\midrule
\textbf{Kf6} & 1 & $0.98598$ & $0.54585$ & $\mathbf{1.53183}$ \\
Kd6 & 1 & $0.96766$ & $0.55696$ & $1.52462$ \\
Kf5 & 0 & $0.66460$ & $0.13279$ & $0.79739$ \\
Kd5 & 0 & $0.66460$ & $0.12983$ & $0.79443$ \\
\bottomrule
\end{tabular}
\end{center}

The two leaders are separated by $0.007$; the two draws have fallen a full $0.7$ behind,
because their FPU value dropped when the good moves got measured and came back healthy.
\textbf{Select Kf6}, and here something new happens: Kf6 already has a child node, so
selection does not stop at the root. It recurses. At the Kf6 node (Black to move, three legal
replies with priors Kf8 $52.05\%$, Kd7 $25.89\%$, Kd8 $22.05\%$), the same rule picks
\textbf{Kf8}, which is expanded and evaluated. The value returned to the root edge Kf6 is
$+0.95129$.

\begin{worked}{following the sign flips on iteration 3}
The new leaf is the position after 1.Kf6 Kf8, \emph{White} to move. Its network value is
$+0.95129$ from White's point of view. Walking up:
\begin{itemize}
\item edge (Kf6-node, Kf8), where \emph{Black} is to move: receives $-0.95129$ --- Black does
  not like it;
\item edge (root, Kf6), where \emph{White} is to move: receives $+0.95129$.
\end{itemize}
Kf6's running average becomes $(0.98598 + 0.95129)/2 = 0.968635$.
\end{worked}

\subsection*{Iterations 4--8 --- the pattern}

The remaining iterations are the same arithmetic; here is the whole run in one table.
"Line" is the path the selection walked before expanding a new node, and "backs up" is the
value that returned to the root edge.

\begin{center}\small
\begin{tabular}{@{}rlrrrrl@{}}
\toprule
It. & Selected & \multicolumn{2}{c}{$S$ of the two leaders} & Line walked & backs up & root edge after \\
\cmidrule(lr){3-4}
 &  & Kd6 & Kf6 & & & \\
\midrule
1 & Kd6 & $1.76358$ & $1.74788$ & Kd6            & $+0.96766$ & Kd6: $n{=}1$, $Q{=}0.96766$ \\
2 & Kf6 & $1.36146$ & $1.52205$ & Kf6            & $+0.98598$ & Kf6: $n{=}1$, $Q{=}0.98598$ \\
3 & Kf6 & $1.52462$ & $1.53183$ & Kf6 Kf8        & $+0.95129$ & Kf6: $n{=}2$, $Q{=}0.96863$ \\
4 & Kd6 & $1.64983$ & $1.41434$ & Kd6 Kd8        & $+0.99759$ & Kd6: $n{=}2$, $Q{=}0.98262$ \\
5 & Kd6 & $1.50779$ & $1.48333$ & Kd6 Kf7        & $+0.99992$ & Kd6: $n{=}3$, $Q{=}0.98839$ \\
6 & Kf6 & $1.42878$ & $1.54411$ & Kf6 Kf8 e6     & $+0.97860$ & Kf6: $n{=}3$, $Q{=}0.97196$ \\
7 & Kd6 & $1.47084$ & $1.44478$ & Kd6 Kd8 e6     & $+0.97060$ & Kd6: $n{=}4$, $Q{=}0.98394$ \\
8 & Kf6 & $1.40085$ & $1.48270$ & \ldots         & & \\
\bottomrule
\end{tabular}
\end{center}

\begin{realdata}[title={The check: hand arithmetic against the engine}]
After seven backed-up values, our pencil-and-paper tree says:
\begin{center}
\begin{tabular}{@{}lrrrr@{}}
\toprule
 & \multicolumn{2}{c}{by hand} & \multicolumn{2}{c}{\texttt{lc0.exe}, \texttt{go nodes 8}} \\
\cmidrule(lr){2-3}\cmidrule(lr){4-5}
Move & $n_a$ & $Q$ & $n_a$ & $Q$ \\
\midrule
Kd6 & 4 & $0.98394$ & 4 & $0.98394$ \\
Kf6 & 3 & $0.97196$ & 3 & $0.97196$ \\
Kf5 & 0 & $0.66651$ & 0 & $0.66652$ \\
Kd5 & 0 & $0.66651$ & 0 & $0.66652$ \\
\bottomrule
\end{tabular}
\end{center}
Every visit count identical; every average identical to five decimal places (the last digit
of the FPU values differs because the priors above are printed to four figures, not because
anything conceptual is missing). \textbf{You have just run Leela's search by hand.}
\end{realdata}

\section{Two things you should notice}

\subsection{The search deepened without being told to}

Nobody wrote "now go deeper". Look at the "line walked" column: iterations 1--2 stop at the
root's children, 3--5 reach two plies, 6--7 reach three. Depth is an \emph{emergent}
property of $U$ collapsing where visits pile up. A node visited $n$ times has its $U$ divided
by $1+n$, so attention flows down the tree the way water flows downhill --- always into the
part of the tree that has been examined least relative to its promise.

Contrast this with alpha--beta, which searches to a fixed depth and then extends
selectively by hand-written rules. Leela has no depth parameter at all. The principal
variation is simply the path you get by always taking the most-visited child.

\subsection{The two bad moves have not been touched}

After eight iterations, Kd5 and Kf5 have $n = 0$. Their FPU value has drifted down from
$0.976$ to $0.667$ while the good moves proved themselves; their $U$ has crept up from
$0.092$ to $0.249$ as the root accumulated visits. They are being ignored, but not
forgotten --- exactly the balance Chapter~\ref{ch:puct} was built to achieve.

So when does the search get to them? Run the engine at higher budgets and watch:

\begin{realdata}[title={Real engine output: the refutation, and its permanence}]
\begin{center}
\begin{tabular}{@{}lrrrrrr@{}}
\toprule
 & \multicolumn{2}{c}{$n$ at 64 nodes} & \multicolumn{2}{c}{$n$ at 128 nodes} & \multicolumn{2}{c}{$n$ at 800 nodes} \\
\cmidrule(lr){2-3}\cmidrule(lr){4-5}\cmidrule(lr){6-7}
Move & $n$ & $Q$ & $n$ & $Q$ & $n$ & $Q$ \\
\midrule
Kd6 & 31 & $0.993$ & 62 & $0.979$ & 377 & $0.956$ \\
Kf6 & 24 & $0.953$ & 58 & $0.963$ & 240 & $0.927$ \\
Kf5 & 0 & --- & \textbf{1} & $\mathbf{0.000}$ & \textbf{1} & $\mathbf{0.000}$ \\
Kd5 & 0 & --- & \textbf{1} & $\mathbf{0.000}$ & \textbf{1} & $\mathbf{0.000}$ \\
\bottomrule
\end{tabular}
\end{center}
Somewhere between 64 and 128 nodes, each of the two drawing moves received \textbf{exactly
one visit}, returned $Q = 0.000$ --- and was never visited again, not in the next 672 nodes.
\end{realdata}

That single visit is the whole story of this chapter in miniature:

\begin{keyidea}
The network's \emph{policy} did not know these moves were bad: it gave them $10.6\%$ of its
attention. The network's \emph{value head} knew immediately: one look returned $0.000$ with a
draw probability of $1.000$. And the \emph{search} is the machinery that converts the second
into a correction of the first, at a cost of two visits out of 800. Policy proposes; value
disposes; search is the negotiation.
\end{keyidea}

And Stockfish, asked independently, agrees that $0.000$ is right: Kd5 and Kf5 draw. Two visits
out of eight hundred were enough to permanently eliminate half the legal moves in the
position, correctly.

\begin{pitfall}
Notice how easily this could have gone wrong, and what it depends on. It worked because the
\emph{value} head recognised the drawn structure in a single forward pass. In a position
where the refutation is five moves deep rather than immediate, that first visit would have
returned an encouraging number, and the search would have needed to build a subtree to find
the truth. The engine's blind spots are precisely the positions where \emph{both} heads are
wrong at once --- and that is not a rare corner case, it is the definition of a difficult
position. Chapter~\ref{ch:readingtree} shows how to spot it happening in your own output.
\end{pitfall}

\section{Which number chooses the move?}

At 800 nodes the search has Kd6 with $n=377$, $Q = 0.956$ and Kf6 with $n = 240$,
$Q = 0.927$. The engine plays \texttt{bestmove e6d6} --- Kd6. Here both criteria agree, so the
position does not settle the question. Which one is it, in general?

The standard answer is \textbf{the most-visited move, not the highest $Q$}, and the reason is
worth understanding because it is counter-intuitive.

\begin{keyidea}
$Q$ is an average over a subtree, and a subtree can be small. A move visited three times with
$Q = 0.99$ has an average over three samples; a move visited 500 times with $Q = 0.95$ has
survived five hundred attempts to refute it. High $Q$ on low $n$ usually means \emph{the
search has not yet found the objection} --- indeed, if the search believed that $Q$, it would
have kept visiting the move, and $n$ would not be small. \textbf{Visit count is the search's
considered opinion; $Q$ is a measurement that may be young.}
\end{keyidea}

The exception is certainty. In Chapter~\ref{ch:puct}'s Morphy position the engine played a
move with $n = 3$ over a move with $n = 172$, because the $Q$ of $1.0$ was not an average at
all --- it was a proof (16.Qb8+ Nxb8 17.Rd8\#). Engines track proven wins and losses
separately from sampled values, and a proof outranks any amount of sampling. That is the only
regular circumstance in which a low-visit move is played.

\begin{repobox}[title={in this repository}]
This is the rule your tooling must follow when reporting "what the engine thought": rank
candidate moves by visits, report $Q$ alongside as the evaluation, and treat any move whose
$Q$ is excellent but whose $n$ is small as \emph{a question}, not an answer. It is also the
right basis for the "how seriously did LC0 weigh this plan" language in your design
documents --- that phrase means visits.
\end{repobox}

\section{An experiment: the position is not the whole input}

One last thing, discovered while checking this chapter's numbers, which matters for any tool
that feeds FENs to a network.

Evaluating the position after 1.Kf6 Kf8 gives different answers depending on how you ask:

\begin{realdata}[title={Real experiment: same board, two histories}]
\begin{center}
\begin{tabular}{@{}lrr@{}}
\toprule
How the position was given to the engine & $V$ & $d$ \\
\midrule
\texttt{position fen 5k2/8/5K2/4P3/8/8/8/8 w - - 2 2} & $+0.98838$ & $0.012$ \\
\texttt{position fen 4k3/\ldots\ w - - 0 1 moves e6f6 e8f8} & $+0.95129$ & $0.049$ \\
\bottomrule
\end{tabular}
\end{center}
The same 32 squares, the same side to move, the same castling and en-passant rights ---
and the evaluation differs by $0.037$, with four times the draw probability. The second
number is the one the search used (it matches the backed-up value in iteration 3 exactly).
\end{realdata}

The cause is that Leela's input is not a position, it is a \textbf{position with history}:
the network sees the last several board states as additional input planes. Given a bare FEN,
the engine has to fabricate a history by repeating the current position, and a repeated
position is evidence of shuffling --- hence the higher draw probability. We measured the same
effect at one ply back in the previous chapter's data ($0.96823$ versus $0.96766$).

\begin{pitfall}
\textbf{This has a direct consequence for your project.} Anywhere the backend evaluates a
position by handing the engine a bare FEN --- diagnosis of a game position, a drill, a cached
critical point --- it is asking a slightly different question from the one the engine answers
during a game. The effect is small in quiet positions and largest where repetition matters:
fortresses, perpetuals, drawn endgames. The fix is cheap: pass the move history
(\texttt{position startpos moves \ldots}) rather than a FEN wherever the history is available.
Worth an audit.
\end{pitfall}

\begin{exercises}

\exhead[warm-up]{ch7:warmup} Reproduce iteration 1's $U$ values for all four moves. Use
$c_{\text{puct}} = 1.7451$ and check all four against the table.

\exhead[the FPU drop]{ch7:fpudrop} Compute $Q_{\text{FPU}}$ at iteration 4, given
$Q(\text{root}) = 0.97024$ and visited prior mass $0.8936$. Why did the visited mass not
change between iterations 3 and 4, even though a visit occurred?

\exhead[iteration 9]{ch7:iter9} Continue the simulation. Iteration 8 selects Kf6 (see the
final row). Suppose it walks the line Kf6 Kf8 e6 Kg8 and backs up $+0.982$. Recompute
$Q(\text{root})$, $Q_{\text{FPU}}$, all four $U$'s, all four $S$'s, and say which move
iteration 9 selects. (Take $c_{\text{puct}} = 1.7459$.)

\exhead[when do the bad moves get looked at?]{ch7:whenbad} At iteration 8 the drawing moves
have $S = 0.915$ against the leader's $1.483$. Using $U = c\,P\sqrt{N}/(1+0)$ with
$P = 0.0538$, and assuming the leaders' $S$ settles near $1.0$ as $U$ decays, estimate the
value of $N$ at which Kf5's $S$ overtakes. Compare with the engine's actual answer (somewhere
between 64 and 128 nodes) and comment on your estimate.

\exhead[the single visit]{ch7:singlevisit} Kf5 was visited once, returned $0.000$, and was
never visited again. Using the PUCT formula, show why: compute $U(\text{Kf5})$ at $N = 800$
with $n = 1$, add $Q = 0$, and compare with the leader's $S$. How large would $N$ have to
become before Kf5 is visited a second time?

\exhead[sign flips]{ch7:signs2} Iteration 6 walked Kf6 Kf8 e6 and the new leaf (Black to
move) evaluated at $-0.97860$. Write out the value added at each of the three edges on the
path, with signs and with whose point of view each represents.

\exhead[policy versus value]{ch7:polvalue} The policy gave the two losing moves $10.6\%$;
the value head scored them $0.000$ on sight. In three sentences, explain why a network can be
wrong in one head and right in the other about the same move. (You will be able to give a
much better answer after Chapter~\ref{ch:heads}; write this one now and compare.)

\exhead[most visits or best Q?]{ch7:visitsvsq} A search reports move A with $n = 410$,
$Q = 0.61$ and move B with $n = 12$, $Q = 0.88$. Which does the engine play, and why? Now
suppose B's $Q$ is exactly $1.000$ and its subtree ends in mate. What changes, and why is
that not an inconsistency?

\exhead[the history effect]{ch7:history} Design an experiment using your own corpus to
measure how large the FEN-versus-history discrepancy is in practice: which positions would
you sample, what would you compare, and what summary statistic would tell you whether it
matters for the diagnosis pipeline? Predict the shape of the answer before you run it.

\exhead[a search that would fail]{ch7:wouldfail} Invent a (schematic, non-chess) situation in
which this search finds the wrong move even with 10{,}000 nodes: specify priors, and the
values returned by the first visit to each move, such that the true best move is never
adequately examined. What single quantity in the PUCT formula would have to change to rescue
it?

\exhead[report writing]{ch7:report} Write the three-sentence explanation your training tool
should show a user about this position, using only quantities computed in this chapter. It
must not contain a chess claim that did not come from the engine.

\end{exercises}
